\documentclass[a4paper,12pt]{article}

\usepackage{mathtools}  
\usepackage{graphicx}

\title{Il mio titolo}

\author{Il mio nome}

\date{\today}

\begin{document}
\maketitle
%
%
\section{La mia sezione}
%
Questa \`e un'equazione con il suo numero
\begin{equation}
S=\int_0^\infty f(x)dx\,.
\label{eq:1}
\end{equation}
Questa equazione pu\`o essere richiama come Eq.~\eqref{eq:1}.
%
\begin{figure}[h]
\begin{center}
\includegraphics[width = 50 mm]{figura.pdf}
\caption{La mia figura.}
\label{fig:1}
\end{center}
\end{figure}\\
%
La fig.~\ref{fig:1} rappresenta un triangolo.
\\
Posso scrivere un'equazione andando a capo
\begin{eqnarray}
\sum_{j=1}^N a_j &=& a_1+a_2+\cdots +a_N\nonumber\\
&=& 45\,.
\label{eq:2}
\end{eqnarray}
%
Possiamo anche scrivere una matrice
$$
A=\left(\begin{array}{cc}
a & b \\
c & d \\
\end{array}\right)\,.
$$
Anche il sistema
\begin{equation*}
\left\{
\begin{array}{rcl}
a_{11}x+a_{12}y&=&b_1\\
a_{21}x+a_{22}y&=&b_2\\
\end{array}\right.\,.
\end{equation*}
Una tabella
\begin{table}[h]
\begin{center}
\begin{tabular}{c|c|c}
Cose & fatti & numeri \\
\hline 
muro & crepa & 3 \\
\hline
\end{tabular}
\caption{Una bella tabella.}
\end{center}
\end{table}
%
\end{document}